% \noindent \textbf{Memory Management}
Recent advancements in scaling large language models (LLMs) to long-horizon tasks have largely focused on overcoming fixed context window constraints through dynamic context management and memory optimization. One primary approach involves context compression and summarization, where methods like ReSum~\cite{wu2025resum}, ACON~\cite{kang2025acon}, and SUPO~\cite{lu2025scaling} employ reinforcement learning to condense interaction histories into dense reasoning states, effectively discarding redundant information while retaining critical evidence. A parallel paradigm explores context folding and Markovian state reconstruction, illustrated by AgentFold~\cite{ye2025agentfold}, Context-Folding~\cite{sun2025scaling}, and The Markovian Thinker~\cite{aghajohari2025markovian}. These frameworks restructure linear history into branching or collapsible sub-trajectories, enforcing Markovian properties to decouple reasoning depth from input length. Finally, approaches such as Mem-$\alpha$~\cite{wang2025mem}, MEM1~\cite{zhou2025mem1}, and DeepMiner~\cite{tang2025beyond} transition from passive context maintenance to active memory construction, training agents via RL to proactively update, retrieve, and synergize external memory stores with reasoning processes, thereby enabling sustained performance over indefinite horizons without linear computational scaling.